\section{Descent to \texorpdfstring{$\GL_n(q)$}{GLn(q)}}
\label{sec:lang-steinberg}

Let \(F:G\to G\) be the standard Frobenius endomorphism, so that
\(G^F=\GL_n(q)\). The chosen representatives are defined over \(\Fq\), hence
their algebraic stabilizers are \(F\)-stable closed subgroups.

\begin{proposition}[Fixed points of the algebraic stabilizer]
\label{prop:fixed-point-stabilizer}
Let \(L\) be an \(F\)-stable representative, and put
\[
  H_L=\Stab_{\GL_n(k)}(L).
\]
Then
\[
  \Stab_{\GL_n(q)}(L)=H_L^F.
\]
\end{proposition}

\begin{proof}
Since \(L\) is defined over \(\Fq\), it is fixed by \(F\). Therefore
\[
  \Stab_{\GL_n(q)}(L)=G^F\cap H_L.
\]
Because \(H_L\) is \(F\)-stable, the right-hand side is exactly \(H_L^F\).
\end{proof}

\begin{theorem}[Finite-field projective quotients]
\label{thm:frobenius-projective-quotients}
Let \(L\) be an \(F\)-stable representative through \(n\leq 7\). Then the
projective quotient
\[
  \overline{Q}_L^F
  =
  \im\bigl(\Stab_{\GL_n(q)}(L)\to \PGL(L)\bigr)
\]
is determined by the Frobenius type of the support divisor \(D(L)\). More
precisely:
\begin{enumerate}[label=\arabic*.]
\item \(D(L)=0\) gives \(\PGL_2(q)\);
\item \(D(L)=m[a]\) with \(a\in \bbP^1(\Fq)\) gives a Borel subgroup;
\item \(D(L)=2[a]+[b]\) with \(a,b\in \bbP^1(\Fq)\) gives a split torus;
\item \(D(L)=[a]+[b]\) with \(a,b\in \bbP^1(\Fq)\) gives the normalizer of a
  split torus;
\item \(D(L)=[\alpha]+[\alpha^q]\) gives the normalizer of a nonsplit torus;
\item \(D(L)=[a]+[b]+[c]\) with \(a,b,c\in \bbP^1(\Fq)\) gives \(S_3\);
\item \(D(L)=[a]+[\alpha]+[\alpha^q]\) gives \(C_2\);
\item \(D(L)=[\alpha]+[\alpha^q]+[\alpha^{q^2}]\) gives \(C_3\).
\end{enumerate}
Consequently the geometric orbit counts \(3,5,11,16\) refine to the finite
field counts \(4,6,14,20\).
\end{theorem}

\begin{proof}
By Theorem~\ref{thm:projective-image}, the algebraic quotient over \(k\) is the
stabilizer of the support divisor \(D(L)\) in \(\PGL_2(k)\). For each divisor
type appearing through \(n\le 7\), the lifting arguments in the proof of
Theorem~\ref{thm:projective-image} are defined over \(\Fq\): diagonal
rescalings, the substitutions \(t\mapsto t/(1+\beta t)\) for one-point support,
transpositions of two simple support summands, and permutations of three
simple support summands. Hence taking \(F\)-fixed points gives exactly the
listed finite quotients.

The only possible Frobenius orbit patterns on the support are split,
quadratic, \(1+2\), and cubic, because the support has size at most \(3\).
Applying these patterns to the geometric families yields the refinements
\(3\mapsto 4\), \(5\mapsto 6\), \(11\mapsto 14\), and \(16\mapsto 20\).
\end{proof}

\begin{lemma}[Frobenius on product kernels]
\label{lem:frobenius-product}
Let \(H_0\) be a connected linear algebraic group over \(k\) with standard
Frobenius \(F_0\). For \(d\geq 1\), let \(H=H_0^d\) and let \(F\) act by
\[
  F(h_1,\dots,h_d)=\bigl(F_0(h_d),F_0(h_1),\dots,F_0(h_{d-1})\bigr).
\]
Then projection onto the first factor induces an isomorphism
\[
  H^F\cong H_0^{F_0^d}.
\]
In particular,
\[
  (\SL_2(k)^d)^F\cong \SL_2(q^d).
\]
\end{lemma}

\begin{proof}
If \((h_1,\dots,h_d)\in H^F\), then
\[
  h_2=F_0(h_1),\quad h_3=F_0^2(h_1),\quad \dots,\quad h_d=F_0^{d-1}(h_1),
\]
and the final relation is \(h_1=F_0^d(h_1)\). So an \(F\)-fixed point is
uniquely determined by its first component, and the displayed isomorphism
follows.
\end{proof}

\begin{theorem}[Finite-field stabilizers from the geometric families]
\label{thm:finite-field-stabilizers}
Let \(L\) be an \(F\)-stable representative of a finite-field orbit through
\(n\leq 7\). Write \(B(q)\), \(T_s(q)\), \(N_s(q)\), and \(N_{ns}(q)\) for a
split Borel subgroup, a split maximal torus, the normalizer of a split maximal
torus, and the normalizer of a nonsplit maximal torus in \(\GL_2(q)\), and let
\(\widetilde{S}_3(q)\), \(\widetilde{C}_2(q)\), and \(\widetilde{C}_3(q)\)
denote the full preimages in \(\GL_2(q)\) of \(S_3\), \(C_2\), and \(C_3\) in
\(\PGL_2(q)\). Then:
\begin{enumerate}[label=\arabic*.]
\item if \(D(L)=0\), then
  \[
    \Stab_{\GL_n(q)}(S_d)
    \cong
    (\Fq^{\,2d-2}\rtimes \Fq^\times)\rtimes \GL_2(q),
  \]
  \[
    \Stab_{\GL_n(q)}(S_2^2)
    \cong
    \bigl((\Sym_2(\Fq)\oplus \Sym_2(\Fq))\rtimes \GL_2(q)\bigr)\rtimes \GL_2(q),
  \]
  and radical extensions are obtained from
  Proposition~\ref{prop:radical-extensions} by fixed points;
\item if \(D(L)=m[a]\) with \(a\in \bbP^1(\Fq)\), then the quotient is
  \(B(q)\), and
  \[
    \Stab_{\GL_n(q)}(J_{a,d})
    \cong
    \SL_2(\Fq[t]/(t^d))\rtimes B(q),
  \]
  \[
    \Stab_{\GL_n(q)}(2J_{a,1})
    \cong
    \Sp_4(q)\rtimes B(q),
  \]
  \[
    \Stab_{\GL_n(q)}(J_{a,2}\oplus J_{a,1})
    \cong
    \bigl(U_{2,1}(q)\rtimes (\SL_2(q)\times \SL_2(q))\bigr)\rtimes B(q),
  \]
  and the families containing an \(S_2\) block are the corresponding
  fixed-point forms of Proposition~\ref{prop:s2-plus-core-kernel} and
  Proposition~\ref{prop:s3-plus-j1};
\item if \(D(L)=[a]+[b]\) with \(a,b\in \bbP^1(\Fq)\), then the quotient is
  \(N_s(q)\), and
  \[
    \Stab_{\GL_n(q)}(J_{a,1}\oplus J_{b,1})
    \cong
    (\SL_2(q)\times \SL_2(q))\rtimes N_s(q);
  \]
\item if \(D(L)=2[a]+[b]\) with \(a,b\in \bbP^1(\Fq)\), then the quotient is
  \(T_s(q)\), and
  \[
    \Stab_{\GL_n(q)}(J_{a,2}\oplus J_{b,1})
    \cong
    (\SL_2(\Fq[t]/(t^2))\times \SL_2(q))\rtimes T_s(q),
  \]
  \[
    \Stab_{\GL_n(q)}(2J_{a,1}\oplus J_{b,1})
    \cong
    (\Sp_4(q)\times \SL_2(q))\rtimes T_s(q);
  \]
\item if \(D(L)=[\alpha]+[\alpha^q]\), then the quotient is \(N_{ns}(q)\), and
  \[
    \Stab_{\GL_n(q)}(Q)\cong \SL_2(q^2)\rtimes N_{ns}(q);
  \]
\item if \(D(L)=[a]+[b]+[c]\) with \(a,b,c\in \bbP^1(\Fq)\), then the quotient
  is \(\widetilde{S}_3(q)\), and
  \[
    \Stab_{\GL_n(q)}(J_{a,1}\oplus J_{b,1}\oplus J_{c,1})
    \cong
    (\SL_2(q)^3)\rtimes \widetilde{S}_3(q);
  \]
\item if \(D(L)=[a]+[\alpha]+[\alpha^q]\), then the quotient is
  \(\widetilde{C}_2(q)\), and
  \[
    \Stab_{\GL_n(q)}(J_{a,1}\oplus Q)
    \cong
    (\SL_2(q)\times \SL_2(q^2))\rtimes \widetilde{C}_2(q);
  \]
\item if \(D(L)=[\alpha]+[\alpha^q]+[\alpha^{q^2}]\), then the quotient is
  \(\widetilde{C}_3(q)\), and
  \[
    \Stab_{\GL_n(q)}(C)\cong \SL_2(q^3)\rtimes \widetilde{C}_3(q).
  \]
\end{enumerate}
In particular, every finite-field stabilizer through \(n\leq 7\) is explicit.
\end{theorem}

\begin{proof}
By Proposition~\ref{prop:fixed-point-stabilizer}, each finite-field stabilizer
is the fixed-point group of the corresponding algebraic stabilizer. The
quotients are given by Theorem~\ref{thm:frobenius-projective-quotients}. For
support defined over \(\Fq\), taking fixed points simply replaces \(k\) by
\(\Fq\) in the formulas of Section~\ref{sec:first-stabilizers}. For quadratic,
\(1+2\), and cubic support, Lemma~\ref{lem:frobenius-product} identifies the
product kernels with \(\SL_2(q^2)\), \(\SL_2(q)\times \SL_2(q^2)\), and
\(\SL_2(q^3)\), respectively. The radical and \(S_2\)-extension families are
again obtained by fixed points, because their defining maps are polynomial in
the matrix entries.
\end{proof}

\begin{center}
\small
\begin{tabular}{L{0.22\textwidth} L{0.22\textwidth} L{0.36\textwidth}}
\toprule
Support pattern over \(\Fq\) & Quotient in \(\GL_2(q)\) & Core kernel over \(\Fq\) \\
\midrule
pure singular &
\(\GL_2(q)\) &
\(\Fq^{\,2d-2}\rtimes \Fq^\times\), or
\((\Sym_2(\Fq)\oplus \Sym_2(\Fq))\rtimes \GL_2(q)\) for \(S_2^2\) \\
one rational support point &
\(B(q)\) &
\(\SL_2(\Fq[t]/(t^d))\), \(\Sp_4(q)\), or
\(U_{2,1}(q)\rtimes (\SL_2(q)\times \SL_2(q))\) \\
two rational points &
\(N_s(q)\) &
\(\SL_2(q)\times \SL_2(q)\) \\
weighted rational pair &
\(T_s(q)\) &
\(\SL_2(\Fq[t]/(t^2))\times \SL_2(q)\), or
\(\Sp_4(q)\times \SL_2(q)\) \\
quadratic pair &
\(N_{ns}(q)\) &
\(\SL_2(q^2)\) \\
three rational points &
\(\widetilde{S}_3(q)\) &
\(\SL_2(q)^3\) \\
one rational point plus one quadratic pair &
\(\widetilde{C}_2(q)\) &
\(\SL_2(q)\times \SL_2(q^2)\) \\
one cubic orbit &
\(\widetilde{C}_3(q)\) &
\(\SL_2(q^3)\) \\
\bottomrule
\end{tabular}
\end{center}

\begin{corollary}[Reconstruction]
\label{cor:finite-field-reconstruction}
Through \(n\leq 7\), every finite-field stabilizer \(\Stab_{\GL_n(q)}(L)\) is
reconstructed from three inputs:
\begin{enumerate}[label=\arabic*.]
\item the geometric support-divisor type \(D(L)\), which determines the
  quotient \(Q_L\);
\item the Frobenius orbit pattern on \(\operatorname{Supp} D(L)\), which
  decides whether the support is split, quadratic, \(1+2\), or cubic;
\item the local repeated-support kernel, together with the radical and
  \(S_2\)-extension mechanisms of Proposition~\ref{prop:radical-extensions},
  Proposition~\ref{prop:s2-plus-core-kernel}, and
  Proposition~\ref{prop:s3-plus-j1}.
\end{enumerate}
\end{corollary}

\begin{proof}
The quotient is determined by
Theorem~\ref{thm:frobenius-projective-quotients}, and the kernel by
Theorem~\ref{thm:finite-field-stabilizers}. This is exactly the reconstruction
procedure.
\end{proof}
